% ============ §15 The dimension of extension rings (原书 pp.116-122) ============
\section{The dimension of extension rings}\label{sec:15}\index{dimension of extension rings, section 15}

\subsection*{1. Fibres}

Let $\varphi\colon A \to B$ be a ring homomorphism, and for $p \in \Spec A$, write $\kappa(p) =
A_p/pA_p$; then $\Spec(B \otimes \kappa(p))$ is called the \emph{fibre}\index{fibre} of $\varphi$ over $p$. As we
saw in \S\ref{sec:7}, it can be identified with the inverse image in $\Spec B$ of $p$ under the
map ${}^a\varphi\colon \Spec B \to \Spec A$ induced by $\varphi$. The ring $B \otimes \kappa(p)$ will be called the
\emph{fibre ring} over $p$. When $(A, m)$ is a local ring, $m$ is the unique closed point of
$\Spec A$, and so the spectrum of $B \otimes \kappa(m) = B/mB$ is called the \emph{closed fibre} of
$\varphi$. If $A$ is an integral domain and $K$ its field of fractions then the spectrum of
$B \otimes_A K = B \otimes_A \kappa(0)$ is called the \emph{generic fibre} of $\varphi$.

\begin{theorem}\label{thm:15.1}
Let $\varphi\colon A \to B$ be a homomorphism of Noetherian rings,
and $P$ a prime ideal of $B$; then setting $p = P \cap A$, we have
\begin{enumerate}[(i)]
  \item $\htop P \le \htop p + \dim B_P/pB_P$;
  \item if $\varphi$ is flat, or more generally if the going-down theorem holds between
        $A$ and $B$, then equality holds in (i).
\end{enumerate}
\end{theorem}

\begin{proof}
We can replace $A$ and $B$ by $A_p$ and $B_P$, and assume that $(A, m)$ and
$(B, n)$ are local rings, with $mB \subset n$. Rewriting (i) in the form
\[
  \dim B \le \dim A + \dim B/mB
\]
makes clear the geometrical content. To prove this, take a system of
parameters $x_1, \ldots, x_r$ of $A$, and choose $y_1, \ldots, y_s \in B$ such that their images
in $B/mB$ form a system of parameters of $B/mB$. Then for $\nu, \rho$ large enough
we have $n^\nu \subset mB + \sum y_iB$ and $m^\rho \subset \sum x_jA$, giving $n^{\nu\rho} \subset \sum y_iB + \sum x_jB$.
Hence $\dim B \le r + s$.

(ii) Let $\dim B/mB = s$, and let $n = P_0 \supset P_1 \supset \cdots \supset P_s$ be a strictly decreas-
ing chain of prime ideals of $B$ between $n$ and $mB$. Obviously we have
$P_i \cap A = m$ for $0 \le i \le s$. Now set $\dim A = r$ and let $m = p_0 \supset p_1 \supset \cdots \supset p_r$
be a strictly decreasing chain of prime ideals of $A$; by the going-down
theorem, we can construct a strictly decreasing chain of prime ideals of $B$
\[
  P_s \supset P_{s+1} \supset \cdots \supset P_{s+r} \quad \text{such that } P_{s+i} \cap A = p_i.
\]
Thus $\dim B \ge r + s$, and putting this together with (i) gives equality.
\end{proof}

\begin{theorem}\label{thm:15.2}
Let $\varphi\colon A \to B$ be a homomorphism of Noetherian rings,
and suppose that the going-up theorem holds between $A$ and $B$. Then if
$p$ and $q$ are prime ideals of $A$ such that $p \supset q$, we have
\[
  \dim B \otimes \kappa(p) \ge \dim B \otimes \kappa(q).
\]
\end{theorem}

\begin{proof}
Set $r = \dim B \otimes \kappa(q)$ and $s = \htop(p/q)$. We choose a strictly
increasing chain $Q_0 \subset Q_1 \subset \cdots \subset Q_r$ of prime ideals of $B$ lying over $q$ and
a strictly increasing chain $q = p_0 \subset p_1 \subset \cdots \subset p_s = p$ of prime ideals of $A$.
By the going-up theorem there exists a chain $Q_r \subset Q_{r+1} \subset \cdots \subset Q_{r+s}$ of
prime ideals of $B$ such that $Q_{r+i} \cap A = p_i$. We set $P = Q_{r+s}$; then
\[
  \htop(P/qB) \ge r + s \quad \text{and} \quad P \cap A = p.
\]
Thus applying the previous theorem to the homomorphism $A/q \to B/qB$
induced by $\varphi$ we get $r + s \le \htop(P/qB) \le s + \dim B_P/pB_P$, and therefore
\[
  r \le \dim B_P/pB_P \le \dim B \otimes \kappa(p). \qquad\blacksquare
\]
\end{proof}

\begin{theorem}\label{thm:15.3}
Let $\varphi\colon A \to B$ be a homomorphism of Noetherian rings, and
suppose that the going-down theorem holds between $A$ and $B$. If $p$ and $q$
are prime ideals of $A$ with $p \supset q$ then
\[
  \dim B \otimes \kappa(p) \le \dim B \otimes \kappa(q).
\]
\end{theorem}

\begin{proof}
We may assume that $\htop(p/q) = 1$, and it is enough to prove that,
given a chain $P_0 \subset P_1 \subset \cdots \subset P_r$ of prime ideals of $B$ lying over $p$ such that
$\htop(P_i/P_{i-1}) = 1$ we can construct a chain of prime ideals $Q_0 \subset Q_1 \subset \cdots \subset Q_r$
of $B$ lying over $q$ such that
\[
  Q_i \subset P_i \quad (0 \le i \le r) \quad \text{and} \quad \htop(Q_i/Q_{i-1}) = 1 \quad (0 < i \le r).
\]
We can find $Q_0$ by going down. If $r \ge 1$ then take $x \in p - q$ and let $T_1, \ldots, T_s$
be the minimal prime divisors of $Q_0 + xB$. Then $\htop(T_j/Q_0) = 1$, while
$\htop(P_r/Q_0) \ge 2$, hence we can choose
\[
  y \in P_r - \bigcup_{j=1}^{s} T_j.
\]
Let $Q_1$ be a minimal prime divisor of $Q_0 + yB$ contained in $P_r$. Then
$\htop(Q_1/Q_0) = 1$, and $y \notin T_i$ for all $i$, hence $x \notin Q_1$.
Therefore $Q_1 \cap A \neq p$, and since $\htop(p/q) = 1$ we must have $Q_1 \cap A = q$. By
the same method we can successively construct $Q_1, Q_2, \ldots, Q_r$.
\end{proof}

\subsection*{2. Polynomial and formal power series rings}

\begin{theorem}\label{thm:15.4}
Let $A$ be a Noetherian ring, and $X_1, \ldots, X_n$ indeterminates
over $A$. Then
\[
  \dim A[X_1, \ldots, X_n] = \dim A[[X_1, \ldots, X_n]] = \dim A + n.
\]
\end{theorem}

\begin{proof}
It is enough to consider the case $n = 1$. For any $p \in \Spec A$, the
fibre $A[X] \otimes_A \kappa(p) = \kappa(p)[X]$ is a principal ideal ring, and therefore one-
dimensional; also $A[X]$ is a free $A$-module, hence faithfully flat, so by
Theorem 1, (ii), $\dim A[X] = \dim A + 1$.

For $A[[X]]$ it is not true in general that $A[[X]] \otimes_A \kappa(p)$ and $\kappa(p)[[X]]$
coincide; however, if $m$ is a maximal ideal of $A$ we have
\[
  A[[X]] \otimes_A \kappa(m) = A[[X]] \otimes_A (A/m) = (A/m)[[X]],
\]
and this fibre ring is one-dimensional. Also, as we saw on p. 4, every
maximal ideal $\mathfrak{M}$ of $A[[X]]$ is of the form $\mathfrak{M} = (m, X)$, where $m = \mathfrak{M} \cap A$
is a maximal ideal of $A$. Thus for a maximal ideal $\mathfrak{M}$ of $A[[X]]$ we have
\[
  \htop \mathfrak{M} = \htop(m \cap A) + 1;
\]
conversely, if $m$ is a maximal ideal of $A$ then $\htop(m, X) = \htop m + 1$, and
putting these together gives $\dim A[[X]] = \dim A + 1$.
\end{proof}

\begin{remark}
It is not necessarily true that a maximal ideal of $A[X]$ lies
over a maximal ideal of $A$. For example, if $A$ is a DVR and $t$ a
uniformising element then $A[t^{-1}] = K$ is the field of fractions of $A$, so
that $A[X]/(tX - 1) \cong K$, and $(tX - 1)$ is a maximal ideal of $A[X]$; however,
$(tX - 1) \cap A = (0)$.
\end{remark}

\begin{remark}
It is quite common for fibre rings of $A \to A[[X_1, \ldots, X_n]]$
to have dimension strictly greater than $n$. For example, let $k$ be a field and
set $A = k[Y, Z]$. It is well-known that the field of fractions of $k[[X]]$ has
infinite transcendence degree over $k(X)$ (see [ZS], vol. II, p. 220). Let $u(X)$,
$v(X) \in k[[X]]$ be two elements algebraically independent over $k(X)$, and
define a $k$-homomorphism (continuous for the $X$-adic topology)
\[
  \varphi\colon A[[X]] = k[Y, Z][[X]] \to k[[X]]
\]
by $\varphi(X) = X$, $\varphi(Y) = u(X)$, $\varphi(Z) = v(X)$. If we set $\Ker \varphi = P$ then $P \cap A = (0)$,
and $A[[X]]/P \cong k[[X]]$ is one-dimensional. Now every maximal ideal of
$A[[X]]$ has height $3$, and, as we will see later, $A[[X]]$ is catenary, so that
$\htop P = 2$. Thus we see that the generic fibre of $A \to A[[X]]$ is two-
dimensional.
\end{remark}

\subsection*{3. The dimension inequality}

We say that a ring $A$ is \emph{universally catenary}\index{universally catenary ring} if $A$ is Noetherian and every
finitely generated $A$-algebra is catenary. Since any $A$-algebra generated
by $n$ elements is a quotient of $A[X_1, \ldots, X_n]$, and since a quotient of a
catenary ring is again catenary, a necessary and sufficient condition for a
Noetherian ring $A$ to be universally catenary is that $A[X_1, \ldots, X_n]$ is
catenary for every $n \ge 0$. (In fact it is known that it is sufficient for $A[X]$
to be catenary, compare Theorem 31.7.)

\begin{theorem}[I. S. Cohen~\cite{Coh3}]\label{thm:15.5}
Let $A$ be a Noetherian integral domain, and
$B$ an extension ring of $A$ which is an integral domain. Let $P \in \Spec B$ and
$p = P \cap A$; then we have
\[
  (*) \quad \htop P + \trdeg_{\kappa(p)} \kappa(P) \le \htop p + \trdeg_K B,
\]
where $\trdeg_K B$ is the transcendence degree of the field of quotients of $B$
over that of $A$.
\end{theorem}

\begin{proof}
We may assume that $B$ is finitely generated over $A$. For if the right
hand side is finite and $m$ and $t$ are non-negative integers such that $m < \htop P$
and $t < \trdeg_{\kappa(p)}\kappa(P)$, then there is a prime ideal chain
$P = P_0 \supset P_1 \supset \cdots \supset P_m$ in $B$. Take $a_i \in P_i - P_{i+1}$, $0 \le i < m$, and let
$c_1, \ldots, c_t \in B$ be such that their images modulo $P$ are algebraically
independent over $A/p$. Set $C = A[\{a_i\}, \{c_j\}]$. If the theorem holds for $C$,
then we have $m + t < \htop p + \trdeg_K C \le \htop p + \trdeg_K B$. Letting $m$ and $t$ vary
we see the validity of $(*)$.

We may furthermore assume, by induction, that $B$ is generated over $A$ by
a single element: $B = A[x]$. We can replace $A$ by $A_p$ and $B$ by $B_P = A_p[x]$,
and hence assume that $A$ is local and $p$ its maximal ideal. Set $k = A/p$ and
write $B = A[X]/Q$. If $Q = (0)$ then $B = A[X]$ and by Theorem 1 we have
$\htop P = \htop p + \htop(P/pB)$, and since $B/pB = k[X]$ we have either $P = pB$ or
$\htop(P/pB) = 1$. In both cases the equality holds in $(*)$.

If $Q \neq (0)$ then $\trdeg_K B = 0$. Since $A$ is a subring of $B$ we have
$Q \cap A = (0)$, so that writing $K$ for the field of fractions of $A$ we have
$\htop Q = \htop QK[X] = 1$. Let $P^*$ be the inverse image of $P$ in $A[X]$. Then
$P = P^*/Q$, $\kappa(P) = \kappa(P^*)$, and
\begin{align*}
  \htop P &\le \htop P^* - \htop Q = \htop P^* - 1 \\
          &= \htop p + 1 - \trdeg_{\kappa(p)}\kappa(P^*) - 1 = \htop p - \trdeg_{\kappa(p)}\kappa(P).
\end{align*}
\end{proof}

\begin{definition}
Suppose that $A$ and $B$ satisfy the conditions of the previous
theorem. We refer to $(*)$ as the \emph{dimension inequality}\index{dimension inequality}, and if the equality in $(*)$
holds for every $P \in \Spec B$, we say that the \emph{dimension formula}\index{dimension formula} holds between
$A$ and $B$. The above proof shows that dimension formula holds between $A$
and $A[x_1, \ldots, x_n]$.
\end{definition}

\begin{theorem}[Ratliff]\label{thm:15.6}
A Noetherian ring $A$ is universally catenary if and
only if the dimension formula holds between $A/p$ and $B$ for every prime
ideal $p$ of $A$ and every finitely generated extension ring $B$ of $A/p$ which
is an integral domain.
\end{theorem}

\begin{proof}[Proof of `only if']
If $A$ is universally catenary then so is $A/p$, so that we
need only consider the case that $A$ is an integral domain, and $B$ is a finitely
generated extension ring which is an integral domain. If
$B = A[X_1, \ldots, X_n]/Q$ and $P = P^*/Q$, then since $A[X_1, \ldots, X_n]$ is catenary
we have $\htop P = \htop P^* - \htop Q$, and an easy calculation proves our assertion.
\end{proof}

\begin{proof}[Proof of `if']
We suppose that $A$ is not universally catenary, so that there
exists a finitely generated $A$-algebra $B$ which is not catenary; without loss
of generality we can assume that $B$ is an integral domain. Write $p$ for the
kernel of the homomorphism $A \to B$. There exist prime ideals $P$ and $Q$
of $B$ such that
\[
  P \subset Q, \quad \htop(Q/P) = d \quad \text{but} \quad \htop Q > \htop P + d.
\]
We write $h = \htop P$, choose $a_1, \ldots, a_h \in P$ such that $\htop(a_1, \ldots, a_h) = h$, and
set $I = (a_1, \ldots, a_h)$, so that $P$ is a minimal prime divisor of $I$. Let
\[
  I = q_1 \cap \cdots \cap q_r
\]
be a shortest primary decomposition of $I$, with $P$ the minimal prime
divisor of $q_1$. Then for $b \in Q \cap q_2 \cap \cdots \cap q_r - P$ we have
\[
  I:b^\nu B = q_1 \quad \text{for } \nu = 1, 2, \ldots.
\]
We set $y_i = a_i/b$ for $1 \le i \le h$,
\[
  C = A[y_1, \ldots, y_h], \quad J = (y_1, \ldots, y_h)C \quad \text{and} \quad M = J + QC = J + Q.
\]
Every element of $C$ can be written in the form $u/b^k$ for suitable $k$, with
$u \in (I + bB)^k$, so that if $z \in J \cap B$ then $zb^\nu \in I$ holds for sufficiently large
$\nu$. Hence $z \in I:b^\nu = q_1$. The converse inclusion $q_1 \subset J \cap B$ is obvious, hence
$J \cap B = q_1$. Thus
\[
  M \cap B = (J + Q) \cap B = (J \cap B) + Q = Q,
\]
\[
  C/J \cong B/q_1, \quad \text{and} \quad C/M \cong B/Q.
\]
Therefore, $C_M/JC_M = B_Q/q_1B_Q$ is a $d$-dimensional local ring, and $J$ is
generated by $h$ elements, so that
\[
  \htop M = \dim C_M \le h + d \le \htop Q.
\]
Now $C$ and $B$ have the same field of fractions, and $\kappa(M) = \kappa(Q)$, so that
this inequality implies that the dimension formula does not hold between
$B$ and $C$. This is a contradiction, since we are assuming that the dimension
formula holds between $A/p$ and $B$ and between $A/p$ and $C$, and one sees
easily that it must then hold between $B$ and $C$.
\end{proof}

\subsection*{4. The Rees ring and gr${}_I$(A)}

Let $A$ be a ring, $I$ an ideal of $A$ and $t$ an indeterminate over $A$. Consider
$A[t]$ as a graded ring in the usual way. We obtain a graded ring $R_+ \subset A[t]$
by setting
\[
  R_+ = R_+(A, I) = \left\{\sum c_nt^n \mid c_n \in I^n\right\} = \bigoplus_{n \ge 0} I^nt^n \subset A[t].
\]
If $I = (a_1, \ldots, a_r)$ then $R_+$ can be written $R_+ = A[a_1t, \ldots, a_rt]$, so that $R_+$ is
Noetherian if $A$ is.

$R_+$ is related to the graded ring $\gr_I(A)$ associated with $A$ and $I$ by the fact
that
\[
  \gr_I(A) = \bigoplus_{n \ge 0} I^n/I^{n+1} \cong R_+/IR_+.
\]

Now let $u = t^{-1}$, and consider $A[t, u] = A[t, t^{-1}]$ as a $\Z$-graded ring in
the obvious way. The \emph{Rees ring}\index{Rees ring} $R(A, I)$ is the graded subring
\[
  R = R(A, I) = R_+[u] = \left\{\sum c_nt^n \in A[t, t^{-1}] \;\middle|\;
    \begin{array}{l}
      c_n \in I^n \text{ for } n < 0 \\
      c_n \in A \text{ for } n > 0
    \end{array} \right\}.
\]
Since
\[
  uR = \left\{\sum c_nt^n \;\middle|\; c_n \in I^{n+1} \text{ for } n \ge 0, \; c_n \in A \text{ for } n < 0\right\},
\]
we have $\gr_I(A) = R/uR$.

\smallskip
\noindent Set $S = \{1, u, u^2, \ldots\}$. Then $R_S = R[u^{-1}] = R[t] = A[t^{-1}, t]$, and
$R_S/(1 - u)R_S = A[t^{-1}, t]/(1 - t) = A$. But $R_S/(1 - u)R_S = (R/(1 - u)R)_{\bar{S}}$,
where $\bar{S}$ is the image of $S$ in $R/(1 - u)R$, and since $\bar{S} = 1$, we see that
$R_{\bar{S}}/(1 - u)R_{\bar{S}} = R/(1 - u)R$. Thus we have
\[
  R/(1 - u)R = A \quad \text{and} \quad R/uR = \gr_I(A),
\]
so that the graded ring $\gr_I(A)$ is a `deformation' of the original ring $A$,
with $R$ as `total space of the deformation', in the sense that $R$ contains a
parameter $u$ such that the values $u = 1$ and $0$ correspond to $A$ and $\gr_I(A)$,
respectively.

We also have
\[
  I^nR \cap A = I^n \quad \text{for all } n \ge 0
\]
and this property is often used to reduce problems about powers of $I$ to
the corresponding problems for powers of the principal ideal $uR$.

We conclude this section by applying the dimension inequality to the
study of the dimension of the Rees ring and $\gr_I(A)$.

Let $A$ be a Noetherian ring, $I = \sum_{i=1}^{r} a_iA$ a proper ideal of $A$, and $t$ an
indeterminate over $A$. We set
\[
  u = t^{-1}, \quad R = R(A, I) = A[u, a_1t, \ldots, a_rt] \quad \text{and} \quad G = \gr_I(A).
\]
We have $R \subset A[t, u]$ and $R/uR \cong G$. For any ideal $\mathfrak{a}$ of $A$, set
\[
  \mathfrak{a}' = \mathfrak{a}A[t, u] \cap R.
\]
Then $\mathfrak{a}' \cap A = \mathfrak{a}A[t, u] \cap A = \mathfrak{a}$, so that for $\mathfrak{a}_1 \neq \mathfrak{a}_2$ we have $\mathfrak{a}_1' \neq \mathfrak{a}_2'$.
Moreover, if $p$ is a prime ideal of $A$ then $p'$ is prime in $R$, and the same
thing goes for primary ideals. If $(0) = q_1 \cap \cdots \cap q_n$ is a primary decompo-
sition of $(0)$ in $A$ then $(0) = q_1' \cap \cdots \cap q_n'$ is a primary decomposition of $(0)$
in $R$. Hence if $p_{0i}$ (for $1 \le i \le m$) are all the minimal prime ideals of $A$ then
$\{p_{0i}'\}_{1 \le i \le m}$ is the set of all minimal prime ideals of $R$. Let $p$ be a prime
ideal of $A$ with $\htop p = h$, and let $p = p_0 \supset p_1 \supset \cdots \supset p_h$ be a strictly decreasing
chain of prime ideals of $A$; then $p' \supset p_1' \supset \cdots \supset p_h'$ is a strictly descending
chain of prime ideals of $R$, so that
\[
  \htop p \le \htop p'.
\]
Conversely, suppose that $P \in \Spec R$ and $P \cap A = p$. Let $p_{0i}$ be a minimal
prime of $R$ contained in $P$ and such that $\htop P = \htop (P/p_{0i}')$; then $R/p_{0i}' \cong
A/p_{0i}$, so that by the dimension inequality
\[
  \htop P = \htop (P/p_{0i}') \le \htop (p/p_{0i}) + 1 - \trdeg_{\kappa(p_{0i})}\kappa(P)
\]
\[
  \le \htop p + 1.
\]
Hence $\dim R \le \dim A + 1$. On the other hand $A[u, t] = R[u^{-1}]$ is a
localisation of $R$ so that $\dim R \ge \dim A[u, t] = \dim A + 1$, so that finally
\[
  \dim R = \dim A + 1.
\]
Moreover, for any $p \in \Spec A$ we set $\bar{a}_i = \text{image of } a_i \bmod p$, so that
\[
  R/p' = (A/p)[u, \bar{a}_1t, \ldots, \bar{a}_rt],
\]
and hence $\trdeg_{\kappa(p)} \kappa(p') = 1$; carrying out the above
calculation using the dimension inequality with $p'$ in place of $P$ we get
$\htop p' \le \htop p$, and so
\[
  \htop p = \htop p'.
\]
We now choose a maximal ideal $m$ of $A$ containing $I$; then since
$R/m' = (A/m)[u]$ we see that $\mathfrak{M} = (m', u)$ is a maximal ideal of $R$ and
$\mathfrak{M} \neq m'$, so that $\htop \mathfrak{M} > \htop m'$. However, by the dimension inequality, we
have $\htop \mathfrak{M} \le \htop m + 1 = \htop m' + 1$. Thus
\[
  \htop \mathfrak{M} = \htop m' + 1 = \htop m + 1.
\]
The element $u$ is a non-zero-divisor of $R$ so that considering a system of
parameters gives $\htop (\mathfrak{M}/uR) = \htop \mathfrak{M} - 1 = \htop m$. Thus providing that there
exists a maximal ideal such that $\htop m = \dim A$ containing $I$, (in particular
if $A$ is local), then we have
\[
  \dim G = \dim (R/uR) = \dim A.
\]
We summarise the above in the following theorem.

\begin{theorem}\label{thm:15.7}
Let $A$ be a Noetherian ring and $I$ a proper ideal; then
setting $R = R(A, I)$ and $G = \gr_I(A)$ we have
\[
  \dim R = \dim A + 1, \quad \dim G \le \dim A.
\]
If in addition $A$ is local, then
\[
  \dim G = \dim A.
\]
\end{theorem}

% ============ 习题 ============
\section*{Exercises to \S\ref{sec:15}}
\emph{Let $k$ be a field.}

\begin{exercise}\label{ex:15.1-ex}
Let $A = k[X, Y] \subset B = k[X, Y, X/Y]$, and $P = (Y, X/Y)B$, $p = (X, Y)A$;
then check that $P \cap A = p$, $\htop P = \htop p = 2$, and $\dim B_P/pB_P = 1$, and hence
that
\[
  \htop P < \htop p + \dim B_P/pB_P.
\]
Show also by a concrete example that the going-down theorem does not
hold between $A$ and $B$.
\end{exercise}

\begin{exercise}\label{ex:15.2-ex}
Does the going-up theorem hold between $A$ and $B$, where $A = k[x] \subset
B = k[x, Y]$?
\end{exercise}

\begin{exercise}\label{ex:15.3-ex}
In Theorem~\ref{thm:15.7}, construct an example where $\dim G < \dim A$.
\end{exercise}
